\section{附录：术语索引与复习路线}

\subsection{术语索引}

\begin{longtable}{p{0.23\textwidth}p{0.32\textwidth}p{0.36\textwidth}}
\toprule
术语 & 一句话解释 & 在本期中的作用 \\
\midrule
Agent & 有边界、在环境中目标导向行动的实体 & 全文核心对象，覆盖人、动物、软件代理和数字世界代理。 \\
Logical Agent & 基于逻辑规则和符号推理的智能体 & 说明 Agent 并非新概念，而是 AI 早期目标。 \\
Semantic Parsing & 把自然语言转成 SQL、API、逻辑式等可执行表示 & Language Agent 的直接前史。 \\
Language Agent & 以 LLM 为语言接口、能规划和使用工具的 Agent & OpenClaw Moment 的技术底座。 \\
ReAct & Reasoning 与 Acting 交替的模式 & 把 LLM 从一次性回答推进到环境交互循环。 \\
Tool Use & 模型调用外部工具/API/函数 & 让模型突破纯文本回答边界。 \\
Computer Use Agent & 能操作浏览器、桌面、移动端或代码环境的 Agent & universal digital agent 的重要入口。 \\
World Model & 对环境状态、因果和可行动性的内部模型 & continual learning 和可靠性的基础。 \\
Continual Learning & 从新交互中持续更新能力 & 2026 年 Agent 发展的主线候选。 \\
Expert Agent & 在某一领域积累技能的专业 Agent & 解决可靠性、速度和成本问题的结果形态。 \\
\bottomrule
\end{longtable}

\subsection{复习路线}

如果只想快速复习本期，可以按四步读：先读第 2 章掌握 Agent 的最小定义；再读第 3--4 章理解技术史和 Language Agent 栈；随后读第 5--7 章理解 OpenClaw、边界消融和 continual learning；最后读第 8--10 章理解大厂 bets、生产评估和社会扩散。这样能把“Agent 为什么突然重要”从产品新闻还原成技术系统问题。

\begin{knowledgebox}{本期与下一期的连接}
EP139 给出 Agent 的技术框架；EP138 罗福莉访谈则会把这个框架放进模型实验室的后训练、RL infra、卡资源分配和组织平权问题里。先读 EP139，再读 EP138，会更容易理解为什么“Agent 范式很吃后训练”。
\end{knowledgebox}

